\documentclass[12pt,a4paper,oneside]{report} % oppure book se vuoi capitoli più "formali"

% --- Encoding e lingua ---
\usepackage[utf8]{inputenc}
\usepackage[T1]{fontenc}
\usepackage[italian]{babel}

% --- Margini ---
\usepackage[a4paper,margin=2.5cm]{geometry}

% --- Matematica ---
\usepackage{amsmath, amssymb, amsthm}

% --- Grafica ---
\usepackage{graphicx}
\usepackage{subcaption}

% --- Linee numerate ---
\usepackage{lineno}

% --- Tabelle ---
\usepackage{booktabs}

% --- Link ---
\usepackage[hidelinks]{hyperref}

% --- Bibliografia ---
\usepackage[backend=biber,style=alphabetic]{biblatex}
\addbibresource{bibliography.bib}

% --- Personalizzazioni ---
\setlength{\parskip}{0.5em}
\setlength{\parindent}{0pt}

% --- Titolo ---
\title{Appunti Tirocinio}
\author{Luca Prigione}
\date{\today}

\begin{document}

\linenumbers % Attiva la numerazione delle linee

\maketitle

% Qui includerai i file separati
\section{Problem Definition}

\subsection{Introduction}

We consider an online variant of the Temporal Hitting Set problem under a sliding window model. 
To this end, we model the input as a temporal hypergraph $H = (V, E)$, where each hyperedge $(S, t)$ consist of a subset of vertices $S \subseteq V$ and a discrete time step $t \in \mathbb{T}$.

Given a time window of size $\epsilon$, we define the set of active hyperedges at time $t$ as:
\[
W(t) = \{ (S, t') \in E \mid t - \epsilon \le t' \le t \}.
\]

At each time step $t$, the algorithm maintains a hitting set $X(t) \subseteq V$ such that:
\[
\forall (S, t') \in W(t), \quad X(t) \cap S \neq \emptyset.
\]

\subsection{Objective}

The goal is to maintain a feasible hitting set over time while optimizing the following criteria:

\begin{itemize}
    \item \textbf{Solution size:} minimize $|X(t)|$ at each time step;
    \item \textbf{Update cost:} minimize the number of changes between consecutive solutions, measured as $\Delta X(t) = |X(t) \setminus X(t-1)| + |X(t-1) \setminus X(t)|$; % Questa formula non è correttissima, perchè guarda solamente il numero di nodi e non effettivamente se i nodi stessi sono uguali o diversi
    \item \textbf{Computational efficiency:} minimize the time required to update the solution as the sliding window evolves.
\end{itemize}

Overall, the objective is to maintain a small and stable hitting set while efficiently handling the insertion and expiration of hyperedges induced by the sliding window.

\end{document}